\documentclass[a4paper,11pt]{article}

\usepackage{fullpage}
\usepackage[all]{xy}
\usepackage[utf8]{inputenc}
\usepackage{amsmath,amsthm}
\usepackage{amsfonts}
\usepackage{graphicx}
\usepackage{latexsym}
\usepackage{float}
\usepackage{color}
%%%%%%%%%%%%%%% COMMANDES %%%%%%%%%%%%%%%%%%%%%%%%%%%

\newcommand{\itemb}{\item[$\bullet$]}

\newcommand{\Fd}{\mathbb{F}}
\newcommand{\Znat}{\mathbb{Z}}
\newcommand{\Nnat}{\mathbb{N}}

%%%%%%%%%%%%%%% ENVIRONEMENTS %%%%%%%%%%%%%%%%%%%%%%%

\theoremstyle{plain}
\newtheorem{lemme}{Lemme}
\newtheorem{algorithme}{Algorithme}
\newtheorem{corolaire}{Corollaire}
\newtheorem{propriete}{Propri\'et\'e}
\newtheorem{proposition}{Proposition}
\newtheorem{theoreme}{Th\'eor\`eme}
\newtheorem{definition}{D\'efinition}

\theoremstyle{definition}
\newtheorem{exercice}{Exercice}
\newtheorem{remarque}{Remarque}
\newtheorem{exemple}{Exemple}


%%%%%%%%%%%%%% DOCUMENTS %%%%%%%%%%%%%%%%%%%%%%%%%%%%


\begin{document}

\hbox{\raisebox{0.4em}{\vrule depth 0pt height 0.5pt width \textwidth}}
\noindent \textbf{Université de Perpignan} \hfill  L3 Informatique \\
 \hfill C. Legrand \hfill  Année \the\year \\
\smallskip
\begin{center}
{
  {\scshape \large TD2 Sécurité \\ Techniques d'infection et polymorphisme}  \\
}
\end{center}
\hbox{\raisebox{0.4em}{\vrule depth 0pt height 0.9pt width \textwidth}}

\bigskip

\bigskip


\begin{exercice}[Bombe logique avec déclencheur temporel]

Une \textit{bombe logique} est un programme malveillant qui s'installe dans un système et attend un événement spécifique appelé \textbf{trigger} (déclencheur) pour exécuter sa fonction offensive (la \textbf{charge utile} ou \textit{payload}).

On considère le programme C suivant :

\begin{verbatim}
#include <stdio.h>
#include <time.h>

int check_trigger() {
    time_t now = time(NULL);
    struct tm *t = localtime(&now);

    // Condition de declenchement
    if (t->tm_mon == 3 && t->tm_mday == 1) {
        return 1;
    }
    return 0;
}

void payload() {
    printf("=== PAYLOAD ACTIVATED ===\n");
    printf("All your files have been encrypted!\n");
}

void normal_operation() {
    printf("System running normally...\n");
}

int main() {
    if (check_trigger()) {
        payload();
    } else {
        normal_operation();
    }
    return 0;
}
\end{verbatim}

\begin{enumerate}
\item Identifier le \textit{trigger} utilisé dans ce programme et expliquer précisément son fonctionnement. À quelle date la charge utile sera-t-elle activée ?

\item Modifier la fonction \texttt{check\_trigger()} pour que le déclenchement se produise uniquement si le programme a été exécuté plus de 100 fois. \textit{Indication : utiliser un fichier pour stocker le compteur d'exécutions.}

\item Proposer deux techniques pour rendre le \textit{trigger} plus difficile à détecter lors d'une analyse statique du code.

\item Du point de vue d'un antivirus, quelles techniques pourraient être utilisées pour détecter ce type de programme malveillant ?
\end{enumerate}

\end{exercice}


\begin{exercice}[Chiffrement XOR et polymorphisme]

Le \textbf{polymorphisme} est une technique utilisée par les virus pour échapper à la détection par signature. En chiffrant le code viral avec des clés différentes à chaque copie, le virus produit des signatures binaires différentes tout en conservant le même comportement.

Le chiffrement \textbf{XOR} (ou exclusif) est souvent utilisé car il possède une propriété intéressante : appliquer deux fois la même opération avec la même clé redonne la valeur originale.

On considère le programme C suivant :

\begin{verbatim}
#include <stdio.h>
#include <string.h>

void xor_crypt(char *data, int len, char key) {
    for (int i = 0; i < len; i++) {
        data[i] = data[i] ^ key;
    }
}

int main() {
    char payload[] = "FORMAT C: /Y";
    char key = 0x42;

    printf("Original : %s\n", payload);

    // Chiffrement
    xor_crypt(payload, strlen(payload), key);
    printf("Chiffre  : ");
    for (int i = 0; i < strlen(payload); i++) {
        printf("%02X ", (unsigned char)payload[i]);
    }
    printf("\n");

    // Dechiffrement
    xor_crypt(payload, strlen(payload), key);
    printf("Dechiffre: %s\n", payload);

    return 0;
}
\end{verbatim}

\begin{enumerate}
\item Expliquer mathématiquement pourquoi l'opération XOR permet à la fois de chiffrer et de déchiffrer avec la même fonction. \textit{Indication : montrer que $A \oplus K \oplus K = A$.}

\item Exécuter mentalement le programme avec la clé \texttt{0x42}. Quel est le premier octet du payload chiffré ? (Le code ASCII de 'F' est 0x46)

\item Modifier la fonction \texttt{xor\_crypt()} pour qu'elle utilise une clé multi-octets (par exemple \texttt{"SECRET"}) au lieu d'un seul octet.

\item Du point de vue d'un antivirus :
\begin{enumerate}
\item Pourquoi cette technique rend-elle la détection par signature inefficace ?
\item Quelle technique alternative l'antivirus pourrait-il utiliser pour détecter ce type de code polymorphe ?
\end{enumerate}
\end{enumerate}

\end{exercice}


\begin{exercice}[Virus compagnon et détournement du PATH]

Un \textbf{virus compagnon} est un type de malware qui ne modifie pas les fichiers existants mais se place à côté d'eux pour être exécuté à leur place. Une technique courante consiste à exploiter la variable d'environnement \texttt{PATH} pour intercepter l'exécution de commandes système.

\textbf{Principe de l'attaque :} Lorsqu'un utilisateur tape une commande (par exemple \texttt{ls}), le système recherche l'exécutable dans les répertoires listés dans \texttt{PATH}, dans l'ordre. Si un attaquant place un faux \texttt{ls} dans un répertoire prioritaire, c'est ce programme malveillant qui sera exécuté.

On considère le programme C suivant, conçu pour remplacer la commande \texttt{ls} :

\begin{verbatim}
#include <stdio.h>
#include <stdlib.h>
#include <unistd.h>

int main(int argc, char *argv[]) {
    // Phase 1: Code malveillant (simulation)
    FILE *f = fopen("/tmp/log.txt", "a");
    if (f) {
        fprintf(f, "Commande interceptee: ls\n");
        fprintf(f, "Repertoire: %s\n", getcwd(NULL, 0));
        fclose(f);
    }

    // Phase 2: Execution de la vraie commande
    execvp("/bin/ls", argv);

    return 0;
}
\end{verbatim}

\begin{enumerate}
\item Expliquer le fonctionnement de ce programme. Pourquoi l'utilisateur ne remarque-t-il pas que son système est compromis ?

\item Comment un attaquant pourrait-il installer ce virus compagnon sur un système Linux ? Proposer deux méthodes différentes.

\item Que se passe-t-il si l'attaquant oublie la ligne \texttt{execvp("/bin/ls", argv)} ? Quelles conséquences cela aurait-il sur la furtivité de l'attaque ?

\item Proposer deux contre-mesures que l'utilisateur ou l'administrateur système pourrait mettre en place pour se protéger contre ce type d'attaque.

\item \textbf{Bonus :} Modifier le programme pour qu'il enregistre également les arguments passés à la commande \texttt{ls}.
\end{enumerate}

\end{exercice}


\begin{exercice}[Routine de recherche de fichiers cibles]

Avant de pouvoir infecter des fichiers, un virus doit d'abord les localiser sur le système. Cette phase de \textbf{recherche de cibles} est cruciale : elle doit être efficace pour maximiser la propagation tout en restant discrète pour éviter la détection.

On considère le programme C suivant qui parcourt récursivement un répertoire à la recherche de fichiers ayant une extension spécifique :

\begin{verbatim}
#include <stdio.h>
#include <stdlib.h>
#include <string.h>
#include <dirent.h>
#include <sys/stat.h>

void search_targets(const char *path, const char *extension) {
    DIR *dir = opendir(path);
    if (dir == NULL) return;

    struct dirent *entry;
    char fullpath[1024];

    while ((entry = readdir(dir)) != NULL) {
        // Ignorer . et ..
        if (strcmp(entry->d_name, ".") == 0 ||
            strcmp(entry->d_name, "..") == 0)
            continue;

        snprintf(fullpath, sizeof(fullpath), "%s/%s", path, entry->d_name);

        struct stat st;
        if (stat(fullpath, &st) == -1) continue;

        if (S_ISDIR(st.st_mode)) {
            // Recursion dans les sous-repertoires
            search_targets(fullpath, extension);
        } else if (S_ISREG(st.st_mode)) {
            // Verifier l'extension
            char *ext = strrchr(entry->d_name, '.');
            if (ext && strcmp(ext, extension) == 0) {
                printf("Cible trouvee: %s (%ld octets)\n",
                       fullpath, st.st_size);
            }
        }
    }
    closedir(dir);
}

int main(int argc, char *argv[]) {
    if (argc != 3) {
        printf("Usage: %s <repertoire> <extension>\n", argv[0]);
        return 1;
    }

    printf("Recherche de fichiers %s dans %s...\n", argv[2], argv[1]);
    search_targets(argv[1], argv[2]);

    return 0;
}
\end{verbatim}

\begin{enumerate}
\item Analyser le code et expliquer le rôle des fonctions \texttt{opendir()}, \texttt{readdir()}, \texttt{stat()} et \texttt{closedir()}.

\item Pourquoi est-il important d'ignorer les entrées "." et ".." lors du parcours d'un répertoire ? Que se passerait-il sinon ?

\item Modifier la fonction \texttt{search\_targets()} pour qu'elle recherche les fichiers exécutables (ayant la permission d'exécution) au lieu de filtrer par extension.
\textit{Indication : utiliser \texttt{st.st\_mode \& S\_IXUSR} pour tester la permission d'exécution.}

\item Le parcours récursif peut être problématique sur de grands systèmes de fichiers. Proposer une modification pour limiter la profondeur de récursion à $N$ niveaux maximum.

\item Du point de vue d'un antivirus, quels comportements suspects ce type de programme pourrait-il déclencher ? Comment un antivirus comportemental pourrait-il détecter cette activité ?
\end{enumerate}

\end{exercice}


\begin{exercice}[Infection par ajout de code (Appending)]

L'\textbf{infection par ajout} (appending) est une technique classique où le virus ajoute son code à la fin d'un fichier cible. Pour éviter d'infecter plusieurs fois le même fichier (ce qui pourrait le corrompre ou le rendre suspect), le virus utilise un \textbf{marqueur d'infection}.

On considère le programme C suivant qui simule ce mécanisme sur des fichiers texte :

\begin{verbatim}
#include <stdio.h>
#include <string.h>
#include <stdlib.h>

#define MARKER "[[INFECTED]]"
#define PAYLOAD "\n--- MALICIOUS CODE SIMULATION ---\n"

int is_infected(const char *filename) {
    FILE *f = fopen(filename, "r");
    if (f == NULL) return -1;

    char buffer[1024];
    while (fgets(buffer, sizeof(buffer), f) != NULL) {
        if (strstr(buffer, MARKER) != NULL) {
            fclose(f);
            return 1;  // Deja infecte
        }
    }
    fclose(f);
    return 0;  // Non infecte
}

int infect_file(const char *filename) {
    if (is_infected(filename) == 1) {
        printf("[SKIP] %s est deja infecte\n", filename);
        return 0;
    }

    FILE *f = fopen(filename, "a");
    if (f == NULL) {
        printf("[ERROR] Impossible d'ouvrir %s\n", filename);
        return -1;
    }

    fprintf(f, "%s", PAYLOAD);
    fprintf(f, "%s\n", MARKER);
    fclose(f);

    printf("[INFECTED] %s\n", filename);
    return 1;
}

int main(int argc, char *argv[]) {
    if (argc < 2) {
        printf("Usage: %s <fichier1> [fichier2] ...\n", argv[0]);
        return 1;
    }

    for (int i = 1; i < argc; i++) {
        infect_file(argv[i]);
    }
    return 0;
}
\end{verbatim}

\begin{enumerate}
\item Expliquer le rôle du marqueur d'infection (\texttt{MARKER}). Que se passerait-il si le virus n'utilisait pas de marqueur ?

\item Pourquoi le fichier est-il ouvert en mode "a" (append) dans la fonction \texttt{infect\_file()} ? Quel serait l'effet d'utiliser le mode "w" à la place ?

\item La fonction \texttt{is\_infected()} parcourt tout le fichier ligne par ligne. Cette approche a des limites en termes de performance. Proposer une méthode plus efficace pour vérifier si un fichier est infecté.
\textit{Indication : penser à la position du marqueur dans le fichier.}

\item Le marqueur \texttt{[[INFECTED]]} est facilement détectable par un antivirus. Proposer une modification pour rendre le marqueur moins visible tout en restant fonctionnel pour le virus.

\item Du point de vue d'un antivirus, quelles caractéristiques de ce programme pourraient être utilisées pour le détecter ? Proposer au moins deux techniques de détection.
\end{enumerate}

\end{exercice}


\end{document}
